\documentstyle[% 


righttag% 


,11pt% 


,amssymb% 


,russian%               remove in Eng. 


,izvru%                 Set izven% in Eng. 


]{amsart} 


 


%


\numberwithin{equation}{section} 


\newcommand{\tts}[1]{} 


 


\theoremstyle{plain} 


\theorembodyfont{\it} 


\newtheorem{propnonum}{\hskip\parindent Предложение} 


\renewcommand{\thepropnonum}{} 


 


%\hoffset=-15mm%                  For Eng. 


\setcounter{page}{11} 


\god{2001} 


\nomer{10~(473)} %                        Remove in Eng. 


%\nomer{Vol.\,45, No.\,10}%               Set in Eng. 


%\str{\pageref{begin}--\pageref{end}}%  Set in Eng. 


\udk{517.938; 513.73 } 


 


\author{\it И.А.\,Галиуллин} 


% NB:   ^^ remove in Eng. 


\title{Бэровский класс показателей Ляпунова механических  


систем, содержащих параметры} 


 


\date{} 


%\pagestyle{myheadings}%         Set in Eng. 


\pagestyle{plain} %              Remove in Eng. 


\begin{document} 


%\thispagestyle{plain}%          Set in Eng. 


\maketitle 


\label{begin} 


 


 


\section{Показатели Ляпунова параметризованных семейств 


дифференциальных уравнений на многообразиях} 


 


Рассматривается связное риманово (бесконечно) 


дифференцируемое многообразие $M^m$ размерности $m$; пусть 


$F$ --- векторное поле на $M^m$, удовлетворяющее условию 


$\sup\limits_{x\in M^m}\|\nabla F(x)\|_1<+\infty$, где $\nabla F$ --- 


ковариантный дифференциал 


векторного поля $F$, норма $\|\cdot\|_1$ --- это операторная норма 


линейного 


отображения, определенная посредством римановой метрики на $M^m$. 


 


Пусть, далее, $S$ --- метрическое пространство (дифференцируемых) 


векторных полей на многообразии $M^m$ и расстояние задано формулой 


$$ 


d(F_1,F_2)=|F_1(x_0)-F_2(x_0)|+\sup_{x\in M^m} 


\|\nabla F_1(x)-\nabla F_2(x)\|_1, 


$$ 


где $x_0$ --- произвольная точка $M^m$, причем выбор другой точки $x_0$ 


приводит к эквивалентной метрике; определенное таким образом 


метрическое пространство, как известно, является полным. Будем 


теперь воспринимать векторное поле $F\in S$ как дифференциальное 


уравнение первого порядка на $M^m$ ([1], с.\,79). 


 


Если $(y_1,\dotsc,y_m)$ --- система локальных координат на открытом 


множестве $U\subset M^m$ и $\sum\limits_i a_i\partial/\partial y_i$ --- 


локальное выражение поля $F$, 


то интегральные кривые $F$ в $U$ являются решениями системы 


дифференциальных уравнений 


\begin{equation} 


\Dot x_i=a_i(x_1,\dotsc,x_m),\quad i=1,\dotsc,m. 


\end{equation} 


 


Как известно, если (1.1) --- система уравнений в $\Bbb R^m$, то система 


уравнений в вариациях (вдоль решения $\wt x(t)$ с начальной точкой 


$\wh x$) имеет вид 


\begin{equation} 


\Dot\xi=A_{F,\Hat x}(t)\xi\quad (A_{F,\Hat x}(t)=F'_x(\wt x(t)), 


\ \; \xi\in\Bbb R^m). 


\end{equation} 


 


Равенства (1.2) представляют линейную систему дифференциальных 


уравнений, для которой известным образом определяются показатели 


Ляпунова $\lambda_1(F,\wh x)\ge\dotsb\ge\lambda_m(F,\wh x)$. 


 


Обобщение понятия $k$-го показателя Ляпунова на случай 


произвольного дифференцируемого риманова многообразия $M^m$ и 


дифференциального уравнения (векторного поля) $F$ на нем проведено 


в [2], а также в [3], где рассматриваются эндоморфизмы $X_t$ 


метризованного векторного расслоения $(E,p_B,B)$ и показатели 


Ляпунова определяются для непрерывного семейства $(t,e)\mapsto|X_t e|$ 


отображений из $\Bbb R^+\times E$ в $\Bbb R$ по формуле 


$$ 


\lambda_k(b)=\inf_{L\in\cal L^{m-k+1}_b} \; 


\underset{t\to+\infty}{\varlimsup}\frac{1}{t}\ln\|X_t|_L\|_2; 


$$ 


здесь $\cal L^{m-k+1}_b$ --- множество всех $(m-k+1)$-мерных векторных 


подпространств слоя $p^{-1}_B(b)$, а $X_t|_L$ --- сужение $X_t$ на $L$, 


норма 


$\|\cdot\|_2$ --- стандартная операторная норма в линейном пространстве 


$L$, 


определенная через риманову метрику метризованного расслоения 


$(E,p_B,B)$; здесь же отмечается, что $\lambda_k$ может принимать значения 


$\pm\infty$ (в частности, $\lambda_k(0)=-\infty$), т.\,е. $\lambda_k$ --- 


отображение $B$ в 


расширенную числовую прямую $\ov{\Bbb R}$. 


 


Пусть $\cal M$ --- некоторое топологическое пространство, элементы 


которого $\mu$ воспринимаются как параметры; пусть $\psi$ --- непрерывное 


отображение $\cal M$ в базу векторного расслоения $(E,p_B,B)$. Для 


случая, когда образ элемента $\mu\in\cal M$ является аргументом показателя 


Ляпунова (что отвечает зависимости элементов базы $B$ от параметра), 


рассматриваются показатели Ляпунова на образе $\psi(\cal M)\subseteq B$, 


зависимость которых от параметров $\mu$ определяется формулой 


$\mu\mapsto\lambda_k(\psi\mu)$ 


(эти функции $\cal M\to\ov{\Bbb R}$ будут обозначаться 


$\ov\lambda_k(\mu)$). Следующее 


утверждение определяет характер этой зависимости и множества 


точек разрыва, когда пространство $\cal M$ метризуемо и полно в некоторой 


метрике. 


 


\begin{thm}[{\series{m}\shape{n}\selectfont [3]}] 


$1$. Для любого $k=1,\dotsc,m$ функция $\ov\lambda_k(\mu):\cal 


M\to\ov{\Bbb R}$ есть 


бэровская функция второго класса $($следовательно, по теореме Бэра 


ее сужение на некоторое всюду плотное множество типа $G_\delta$ в $\cal 


M$ является непрерывной функцией$)$. 


 


$2$. Точка полунепрерывности сверху каждой из функций 


$\ov\lambda_k(\mu):\cal M\to\ov{\Bbb R}$ 


является типичной в $\cal M$ $($это означает, что множество 


таких точек является в $\cal M$ всюду плотным типа $G_\delta)$. 


\end{thm} 


 


Пусть $\Dot x=F(x,\mu)$ --- 


дифференциальное уравнение ``с правой частью, зависящей 


от параметров''. Конструкции работы [3] таковы, что если потребовать 


непрерывную зависимость элементов базы, т.\,е. поля $F$ и точек 


многообразия (``начальных значений''), от параметров, то применимость 


приведенной теоремы к рассматриваемой ситуации, когда векторное 


поле либо структура многообразия или и то, и другое (непрерывно) 


зависят от параметра, будет вытекать из следующего утверждения, 


являющегося частным вариантом теоремы ([4], с.\,385) о бэровском классе 


сложной функции, когда известны классы обеих функций. 


 


\begin{propnonum} 


Пусть $\cal M_0$ и $B_0$ --- метрические пространства, 


отображение $\psi_0:\cal M_0\to B_0$ непрерывно и функция 


$\lambda_0(\cdot):\psi_0(\cal M_0)\to\ov{\Bbb R}$ 


$(\ov{\Bbb R}$ --- расширенная числовая прямая$)$ принадлежит второму 


бэровскому 


классу. Тогда функция 


$\ov\lambda_0(\cdot)\equiv\lambda_0(\psi_0(\cdot)):\cal M_0\to\ov{\Bbb R}$ 


есть также 


функция второго бэровского класса. 


\end{propnonum} 


 


К рассмотренному случаю приложимы все утверждения теоремы 1, 


когда вместо уравнения (1.1) исследуется неавтономная система, 


т.\,е. уравнение вида $\Dot x=F(x,t,\mu)$, 


где векторное поле $F$ непрерывно зависит от параметра $\mu\in\cal M$. 


 


\section{Реономные механические системы и дифференциальные 


уравнения второго порядка на многообразиях} 


 


Здесь обобщаются необходимые конструкции [1], на основе которых 


устанавливается возможность применения теоремы 1 к классу 


векторных полей $\cal F$ на касательном расслоении, определяемых как 


дифференциальные уравнения второго порядка. 


 


Пусть $T(M)$ --- касательное пространство многообразия $M^m$ без 


края, которое также является многообразием, причем, как следует 


из структуры его атласа ([1], с.\,68), оно будет связным, если таковым 


является само многообразие $M^m$. Локальное выражение поля $\cal F$ 


имеет вид ([1], с.\,159) 


\begin{equation} 


\sum_i\bigg(\Dot q_i\frac{\partial}{\partial q_i}+a_i


(q_1,\dotsc,q_m;\,\Dot q_1,\dotsc,\Dot q_m) 


\frac{\partial}{\partial\Dot q_i}\bigg). 


\end{equation} 


 


Интегральные кривые поля $\cal F$ локально являются решениями 


системы дифференциальных уравнений 


\begin{equation} 


\frac{dq_i}{dt}=\Dot q_i,\quad 


\frac{d\Dot q_i}{dt}=a_i(q_1,\dotsc,q_m;\, 


\Dot q_1,\dotsc,\Dot q_m),\quad i=1,\dotsc,m. 


\end{equation} 


 


Изложенные в ([1], с.\,152--170) определения и свойства далее 


используются для обобщения понятия механической системы


([1], с.\,173), где конфигурационное пространство определяется как 


многообразие $M^m$, фазовое пространство --- как его касательное 


пространство $T(M)$, кинетическая энергия --- это некоторая 


дифференцируемая функция $T$ на $T(M)$, а с силовым полем 


связана полубазисная форма Пфаффа $\pi$ на $T(M)$ ([1], с.\,167), наконец, 


сама механическая система определена как тройка $(M^m,T,\pi)$. 


 


Список аргументов функции $T$ в локальном выражении, таким 


образом, состоит из ``обобщенных координат'' $q_1,\dotsc,q_m$ и 


``обобщенных 


скоростей'' $\Dot q_1,\dotsc,\Dot q_m$, что соответствует классу 


механических систем, 


которые называются склерономными ([5], с.\,35--36). Ниже приводится 


построение математических объектов, позволяющих на языке 


дифференциальной геометрии описать реономные механические 


системы, т.\,е. такие, для которых кинетическая энергия $T$ зависит 


также и от времени. В этом случае $T$ определяется здесь как 


дифференцируемая функция на многообразии $T(M)\times\Bbb R$, ее локальное 


выражение имеет вид $T=T(q_1,\dotsc,q_m;\,\Dot q_1,\dotsc,\Dot q_m;t)$. 


На самом 


касательном пространстве $T(M)$ полагаются определенными два 


дифференциальных оператора ([1], с.\,164--165), представляющих собой 


эндоморфизмы пространства дифференциальных форм 


$\Lambda(T(M))$ на $T(M)$: 


вертикальное дифференцирование $\bold i_v$, действующее по формуле 


$\bold i_v\alpha(\cal X_1,\dotsc,\cal X_p)\mapsto \sum\limits_i \alpha(\cal 


X_1,\dotsc,v\cal X_i,\dotsc,\cal X_p)$, 


когда $\alpha$ --- форма степени $p$, и $\bold i_v f=0$ для дифференцируемой 


функции $f$ на $T(M)$ ($\cal X_i$ --- векторное поле, a $v$ --- 


вертикальный 


эндоморфизм второго касательного расслоения ([1], с.\,161)); 


вертикальное антидифференцирование $\bold d_v$ по формуле


$\bold d_v=\bold i_v \bold d-\bold d \bold i_v$, 


где $\bold d$ --- внешнее дифференцирование. 


 


В качестве обобщений этих определений здесь вводятся следующие 


операторы на $\Lambda(T(M)\times\Bbb R)$, где $T(M)\times\Bbb R$ --- 


дифференцируемое 


многообразие размерности $2m+1$. 


 


\begin{defn} 


Вертикальным варьированием $\Lambda(T(M)\times\Bbb R)$ 


называется оператор $\boldsymbol \iota_v$, при каждом $t\in\Bbb R$


являющийся вертикальным 


дифференцированием $\Lambda(T(M))\subset\Lambda(T(M)\times\Bbb R)$. 


\end{defn} 


 


\begin{defn} 


Вертикальным антиварьированием $\Lambda(T(M)\times\Bbb R)$ 


называется оператор $\boldsymbol \delta_v$, при каждом $t\in\Bbb R$


являющийся вертикальным 


антидифференцированием $\Lambda(T(M))\subset\Lambda(T(M)\times\Bbb R)$. 


\end{defn} 


 


Тогда имеют место локальные формулы, аналогичные ([1], с.\,164), 


\begin{gather*} 


\boldsymbol \iota_v f=0,\quad \boldsymbol \iota_v(dq_i)=


\boldsymbol \iota_v(dt)=0,\quad 


\boldsymbol \iota_v(d\Dot q_i)=dq_i,\\ 


\boldsymbol \delta_v f=\sum_i\frac{\partial f}{\partial\Dot q_i}dq_i,\quad 


\boldsymbol \delta_v(dq_i)=\boldsymbol \delta_v(d\Dot q_i)=


\boldsymbol \delta_v(dt)=0,\quad 


i=1,\dotsc,m. 


\end{gather*} 


 


\begin{defn} 


Внешним варьированием на многообразии $M^m\times\Bbb R$ 


называется оператор $\boldsymbol \delta$, при каждом $t\in\Bbb R$


совпадающий с внешним 


дифференцированием $\Lambda(M^m)\subset \Lambda(M^m\times\Bbb R)$. 


\end{defn} 


 


Аналогично ([1], с.\,89) оператор $\boldsymbol \delta$ локально


определяется формулами 


$\boldsymbol \delta f=d(f\big|_{t=\text{const}})$


и $\boldsymbol \delta(d f)=0$, где $f$ --- 


дифференцируемая функция 


на $M^m\times\Bbb R$. 


 


Очевидно, операторы $\boldsymbol \iota_v$


и $\boldsymbol \delta_v$ (соответственно $\boldsymbol \delta$) 


наследуют 


свойства операторов $\bold i_v$ и $\bold d_v$


на $T(M)$ (соответственно $\bold d$ на $M^m$). 


 


Как отмечалось выше, $T(M)\times\Bbb R$ является многообразием, и на 


нем среди всех векторных полей будут выделяться те поля, которые 


имеют вид $\cal X+\partial/\partial t$, где $\cal X$ --- векторное поле на 


$T(M)$. 


 


\begin{defn} 


Неавтономным дифференциальным уравнением 


второго порядка $\cal T$ называется векторное поле $\cal 


F+\partial/\partial t$ на $T(M)\times\Bbb R$, 


для которого поле $\cal F$ является дифференциальным уравнением 


второго порядка на $T(M)$ ([1], с.\,159). 


 


При этом само поле $\cal F$ локально имеет выражение (2.1). 


\end{defn} 


 


\begin{defn}[{\series{m}\shape{n}\selectfont ср. [1], c.\,171}] 


Замкнутая форма $\omega=\bold d\boldsymbol \delta_v T$ 


степени 2 на $T(M)\times\Bbb R$ называется фундаментальной формой 


реономной механической системы $(M^m,T,\pi)$. 


\end{defn} 


 


Ниже будет предполагаться, что сужение $\omega|_{T(M)}$ --- 


симплектическая 


форма на $T(M)\times\{t\}$, тогда многообразие $T(M)$ является 


симплектическим. 


 


Пусть $\bold i_{\cal Y}$ --- оператор внутреннего произведения на некоторое 


векторное поле $\cal Y$. 


 


\begin{lem} 


Дифференциальная форма


$\bold i_{\cal X+\partial/\partial t}(\bold d\boldsymbol \delta_v T)$ 


является 


формой Пфаффа, принадлежащей $\Lambda(T(M)\times\Bbb R)$. 


\end{lem} 


 


{\bf Доказательство.} Достаточно использовать локальные выражения. 


Так как 


$$ 


\boldsymbol \delta_v T=\sum_i\frac{\partial T}{\partial \Dot q_j}dq_j, 


$$ 


то 


$$ 


\omega=\bold d\boldsymbol \delta_v T=\sum_{i,j} 


\frac{\partial^2T}{\partial\Dot q_i\partial\Dot q_j}d\Dot q_i 


\wedge dq_j+\sum_{i,j} 


\frac{\partial^2T}{\partial q_i\partial\Dot q_j}dq_i\wedge 


dq_j+\sum_j 


\frac{\partial^2T}{\partial t\partial\Dot q_j}dt\wedge dq_j.


$$ 


 


Пусть векторное поле $\cal X$ на $T(M)$ выражается формулой 


\begin{equation} 


\sum_i\bigg(a_i(q_1,\dotsc,q_m;\,\Dot q_1,\dotsc,\Dot q_m) 


\frac{\partial}{\partial\Dot q_i}+b_i(q_1,\dotsc,q_m;\; 


\Dot q_1,\dotsc,\Dot q_m)\frac{\partial}{\partial q_i}\bigg). 


\end{equation} 


Внутреннее произведение $\bold i_{\cal Y}$ действует по правилам:


$\bold i_{\cal 


Y}(df)=\cal Y f$ 


для дифференцируемой функции $f$; $\bold i_{\cal Y}(\alpha\wedge\beta)= 


(\bold i_{\cal Y}(\alpha))\wedge\beta-\alpha\wedge(\bold i_{\cal Y}(\beta))$ 


для дифференциальных форм $\alpha$ и $\beta$. В частности, $\bold i_{\cal 


X}(d\Dot q_i)=a_i$, 


$\bold i_{\cal X}(dq_i)=b_i$, $\bold i_{\cal X}(dt)=0$, 


$\bold i_{\partial/\partial t}(dq_i)=\bold


i_{\partial/\partial t}(d\Dot q_i)=0$, 


$\bold i_{\partial/\partial t}(dt)=1$. 


Поэтому 


\begin{multline} 


\bold i_{\cal X+\partial/\partial t}\omega=\bold i_{\cal X}\omega+ 


\bold i_{\partial/\partial t}\omega=\sum_j\bigg(\sum_i 


\frac{\partial^2T}{\partial\Dot q_i\partial\Dot q_j}a_i+ 


\frac{\partial^2T}{\partial t\partial\Dot q_j}\bigg)dq_j- \\ 


% 


-\sum_{i,j}\frac{\partial^2T}{\partial\Dot q_i\partial\Dot q_j} 


b_i\,d\Dot q_j-\sum_i 


\frac{\partial^2T}{\partial t\partial\Dot q_i}b_i\,dt.\quad\square 


\end{multline} 


 


Среди полей на $T(M)$ общего вида (2.3) особо выделяется векторное 


поле Лиувилля $\cal V$, порождающее так называемую однопараметрическую 


группу гомотетий ([1], с.\,158), которые локально задаются формулой 


$(q_i,\Dot q_i)\mapsto(q_i,e^t\Dot q_i)$, т.\,е. для $\cal V$ 


коэффициенты в (2.3) имеют вид 


$a_i=\Dot q_i$, $b_i=0$, $i=1,\dotsc,m$. 


 


Из определения 3 теперь выводится


 


\begin{lem} 


Пусть $(M^m,T,\pi)$ --- реономная механическая система. 


Тогда существует и притом единственное векторное поле $\cal Z$ на 


$T(M)\times\Bbb R$ такое, что 


\begin{equation} 


\bold i_{\cal Z}(\bold d\boldsymbol \delta_v T)=


\boldsymbol \delta T-d(\cal V T)+\pi, 


\end{equation} 


где $\cal V$ --- векторное поле Лиувилля на $T(M)$. 


\end{lem} 


 


\begin{pf} 


Согласно утверждению ([1], с.\,126) отображение 


$\cal X\mapsto \bold i_{\cal X}\omega$ выражает изоморфизм векторных полей 


$\cal X$ на 


симплектическом многообразии с одной стороны и дифференциальных 


форм Пфаффа на этом многообразии --- с другой. Правая часть 


равенства (2.5), очевидно, является линейной дифференциальной 


формой, элементом пространства $\Lambda(T(M)\times\Bbb R)$. По лемме 1 


аналогичное 


справедливо в отношении левой части. Зафиксируем произвольное 


$t\in\Bbb R$. При 


$\cal Z=\cal X+\frac{\partial}{\partial t}$, 


где $\cal X$ --- векторное поле на $T(M)$, (2.5) 


превращается в равенство $\bold i_{\cal X}(d d_v T)=d(T-\cal V T)+\pi$, 


являющееся 


следствием указанного изоморфизма. Следовательно, равенство 


(2.5) имеет место на прямом произведении $\Lambda(T(M)\times\Bbb R)$, и из 


существования и единственности поля $\cal X$ вытекает существование 


и единственность поля 


$\cal Z=\cal X+\frac{\partial}{\partial t}$.


\end{pf}


 


Аналогично определению ([1], с.\,174) векторное поле $\cal Z$ будет 


называться динамической системой, соответствующей (реономной) 


механической системе $(M^m,T,\pi)$. Следующая теорема конкретизирует 


связь между рассматриваемыми полями и системами. 


 


\begin{thm} 


Динамическая система $\cal Z$, соответствующая 


реономной механической системе $(M^m,T,\pi)$, является неавтономным 


дифференциальным уравнением второго порядка $\cal T$ на многообразии 


$M^m$. 


\end{thm} 


 


\begin{pf} 


Используя определение поля Лиувилля $\cal V$ и 


оператора $\boldsymbol \delta$, можно записать 


$$ 


\cal V T=\sum_i\frac{\partial T}{\partial\Dot q_i}\Dot q_i;\quad 


\boldsymbol \delta T=\sum_j\bigg( 


\frac{\partial T}{\partial q_j}d q_j+ 


\frac{\partial T}{\partial\Dot q_j}d\Dot q_j\bigg). 


$$ 


Пусть, далее, локальное выражение для $\pi$ имеет вид 


$\pi=\sum\limits_j Q_j d q_j$. 


Тогда правая часть равенства (2.5) представляет дифференциальную 


форму 


\begin{equation} 


\boldsymbol \delta T-d(\cal V T)+\pi=\sum_j\bigg( 


\frac{\partial T}{\partial q_j}-\sum_i 


\frac{\partial^2T}{\partial q_j\partial\Dot q_i}\Dot q_i+Q_j\bigg) 


d q_j-\sum_{i,j} 


\frac{\partial^2T}{\partial\Dot q_i\partial\Dot q_j} 


\Dot q_i d\Dot q_j-\sum_i 


\frac{\partial^2T}{\partial t\partial\Dot q_i}\Dot q_i dt. 


\end{equation} 


Из сравнения коэффициентов при $d\Dot q_j$ в выражениях (2.4) и (2.6) 


следуют равенства $b_i=\Dot q_i$ ($i=1,\dotsc,m$), что приводит к формуле 


\begin{equation} 


\cal Z=\sum_i\bigg(\Dot q_i 


\frac{\partial}{\partial q_i}+a_i 


\frac{\partial}{\partial\Dot q_i}\bigg)+ 


\frac{\partial}{\partial t}. 


\end{equation} 


 


В силу (2.1) это означает, что векторное поле $\cal Z\equiv\cal T$ 


является 


(неавтономным) дифференциальным уравнением второго порядка. 


\end{pf} 


 


Аналогично, сравнение коэффициентов при $d q_j$ с учетом уравнений 


(2.2) дает систему равенств 


$$ 


\sum_i\frac{\partial^2 T}{\partial\Dot q_i\partial\Dot q_j} 


\frac{d\Dot q_i}{dt}+\sum_i 


\frac{\partial^2T}{\partial q_i\partial\Dot q_j} 


\frac{d q_i}{dt}+ 


\frac{\partial^2T}{\partial t\partial\Dot q_j}= 


\frac{\partial T}{\partial q_j}+Q_j,\quad j=1,\dotsc,m, 


$$ 


что эквивалентно системе дифференциальных уравнений, имеющих 


форму классических уравнений Лагранжа 


$$ 


\frac{d}{dt} 


\bigg(\frac{\partial T}{\partial\Dot q_j} \bigg)


-\bigg(\frac{\partial T}{\partial q_j}\bigg)=Q_j,\quad 


j=1,\dotsc,m,


$$ 


и описывающих движение как склерономных, так и реономных 


механических систем. 


 


\section{Показатели Ляпунова голономных механических систем} 


 


Здесь устанавливается бэровский класс показателей Ляпунова 


эндоморфизмов, определяемых векторными полями особого рода, а 


именно, (неавтономными) дифференциальными уравнениями второго 


порядка. Исходным объектом является механическая система в 


смысле определения предыдущего параграфа, притом содержащая 


параметры, более точно, непрерывно зависящая от них. Последнее 


означает, что выполняются условия теоремы 1 и вытекающие из 


нее утверждения. Указанное векторное поле --- это динамическая 


система $\cal T$, соответствующая (реономной) механической системе 


$(M^m,T,\pi)$. Векторное расслоение в этой ситуации имеет структуру, 


описанную в \S\,1. Итак, предполагается, что существует непрерывное 


отображение $\psi$ полного метрического пространства $\cal M$ (параметров) 


в базу $B$ расслоения, одним из составляющих элементов в которой 


является совокупность векторных полей на многообразии $T(M)\times\Bbb R$. 


 


В \S\,2 от многообразия $M^m$ не требовалось, чтобы оно было 


римановым, здесь это требование существенно, причем многообразие 


$T(M)$ также можно считать римановым (в отдельных случаях 


метрикой служит кинетическая энергия $T$). Риманову структуру 


можно задать и на следующем касательном расслоении, для которого 


тотальное пространство $T(T(M))$ называется вторым касательным 


пространством ([1], с.\,157). 


 


Теорема 1, изложенная в \S\,1, обладая высокой степенью общности, 


предусматривает, таким образом, проверку выполнения единственного 


условия, а именно, неравенства (1.1) в форме 


\begin{equation} 


\sup_{x\in T(M)}\|\nabla\cal T(x)\|_3<+\infty, 


\end{equation} 


где норма $\|\cdot\|_3$ аналогична норме $\|\cdot\|_1$. 


 


Пусть векторное поле $\cal T$ в локальных координатах имеет 


вид 


(2.7), где векторы $\partial/\partial q_i$, $\partial/\partial \Dot q_i$, 


$\partial/\partial t$, $i=1,\dotsc,m$, являются 


базисными в касательном расслоении многообразия $T(M)\times\Bbb R$. 


Определение нормы ковариантного дифференциала требует 


рассмотрения ковариантной производной $\nabla_{\cal Y}\cdot\cal T$ 


относительно 


некоторого векторного поля $\cal Y$ на этом многообразии, которое также 


раскладывается по этим векторам. Из линейности оператора по 


аргументу $\cal Y$ следует, что для исследования неравенства (3.1) 


достаточно в качестве $\cal Y$ вначале выбрать любой из базисных векторов. 


Тогда по формулам ([6], сс.\,79, 82) (в частности, для 


дифференцируемой функции 


$f:\nabla_{\cal Y}\cdot(f\cal T)=f\nabla_{\cal Y}\cdot\cal F+\cal Y 


f\cdot \cal F$) 


получаются следующие выражения для ковариантных производных: 


\begin{multline} 


\nabla_{\cal Y}\cal T=\sum^m_{i=1}\bigg[\Dot q_i\bigg( 


\sum^m_{\nu=1}\Gamma^\nu_{i,(\cdot)} 


\frac{\partial}{\partial q_\nu}+\sum^{2m}_{\nu=m+1} 


\Gamma^\nu_{i,(\cdot)}\frac{\partial}{\partial\Dot q_\nu} 


+\Gamma^{2m+1}_{i,(\cdot)}\frac{\partial}{\partial t}\bigg)+ \\ 


% 


+a_i\bigg(\sum^m_{\nu=1}\Gamma'{}^\nu_{i,(\cdot)} 


\frac{\partial}{\partial q_\nu}+\sum^{2m}_{\nu=m+1} 


\Gamma'{}^\nu_{i,(\cdot)}\frac{\partial}{\partial\Dot q_\nu} 


+\Gamma'{}^{2m+1}_{i,(\cdot)}\frac{\partial}{\partial t}\bigg)+ 


\frac{\partial a_i}{\partial(\ast)}\cal T\bigg]+\\ 


% 


+\sum^m_{\nu=1}\Gamma^\nu_{2m+1,(\cdot)} 


\frac{\partial}{\partial q_\nu}+\sum^{2m}_{\nu=m+1} 


\Gamma^\nu_{2m+1,(\cdot)}\frac{\partial}{\partial\Dot q_\nu}+ 


\Gamma^{2m+1}_{2m+1,(\cdot)}\frac{\partial}{\partial t} 


+\cal T, 


\end{multline} 


где функции $\Gamma^\nu_{i,j}$ и $\Gamma'{}^\nu_{i,j}$ --- компоненты, или 


символы Кристоффеля 


линейной связности, через $(\cdot)$ обозначен номер базисного вектора 


$\cal Y$, 


через $\partial/\partial(\ast)$ --- производная по соответствующей 


координате. 


 


 


В этой формуле следует вместо $\cal T$ написать выражение (2.7) и 


собрать коэффициенты при $\partial/\partial q_i$, 


$\partial/\partial\Dot q_i$, $\partial/\partial t$, $i=1,\dotsc,m$. 


Теперь условие (3.1) оказывается эквивалентным ограниченности 


этих коэффициентов, когда $\cal Y$ пробегает множество всех базисных 


векторов, благодаря тому, что в конечномерном пространстве норму 


можно определить как максимальный коэффициент растяжения 


векторов ортонормированного базиса. 


 


Далее, т.\,к. понятие механической системы приобрело свое 


формализованное определение, стало возможным говорить о 


показателях Ляпунова механической системы (реономной или 


склерономной), имея в виду показатели Ляпунова систем уравнений 


в вариациях, т.\,е. функции $\lambda_k(b)$ или же функции 


$\ov\lambda_k(\mu)$, $k=1,\dotsc,2m$, когда 


система содержит параметры. 


 


Как известно, если в некотором многообразии $M^m$ выделить 


открытое подмножество $V$, то дифференциальная структура $M^m$ 


индуцирует на $V$ структуру дифференцируемого многообразия, 


которое называется открытым подмногообразием ([1], с.\,56). 


 


Структура уравнений Лагранжа на каждой карте такова ([7], с.\,52--54), 


что из-за наличия соответствующих линейных и квадратичных 


по $q$ и $\Dot q$ членов условие ограниченности в разложении (3.2) в 


большинстве случаев может выполняться лишь на некотором 


открытом подмногообразии многообразия $T(M)$. Тогда в качестве 


указанного отображения $\psi$ выбирается его ограничение 


на соответствующем прообразе. В частности, для параллелизуемого 


многообразия условие ограниченности коэффициентов в (3.2) 


достаточно проверять на $M^m$, а в $\Bbb R^m$ брать любой шар достаточно 


большого радиуса; для многообразия с атласом, состоящим из 


конечного числа карт, из всех радиусов выбирается наибольший. 


 


Содержание этого параграфа составляет теперь доказательство 


следующего утверждения, являющегося результатом применения 


теоремы 1 к рассматриваемому классу дифференциальных 


уравнений на многообразиях. 


 


\begin{thm} 


Пусть на некотором открытом подмногообразии 


пространства, касательного к конфигурационному, выполнены 


условия ограниченности коэффициентов в разложениях $(3.2)$. 


Тогда показатели Ляпунова механической системы, содержащей 


параметры, как функции, определенные на полном метрическом 


пространстве параметров, принадлежат второму бэровскому классу. 


\end{thm} 


 


Из этой теоремы выводятся как следствие другие утверждения, 


сформулированные в теореме 1, в частности, о том, что в пространстве 


параметров для каждого показателя Ляпунова точка 


полунепрерывности сверху является типичной. 


 


Автор выражает глубокую благодарность В.М.\,Миллионщикову 


за консультации. 


 


\begin{thebibliography}{9} 


 


\bibitem{1} 


Годбийон К. {\em Дифференциальная геометрия и аналитическая 


механика}. -- М.: Мир, 1973. -- 188 с. 


\bibitem{2} 


Миллионщиков В.М. {\em Бэровские классы функций и 


показатели Ляпунова}. III // Дифференц. уравнения. -- 1980. 


-- Т.\,16. -- \No\,10. -- С.\,1766--1785. 


\bibitem{3} 


Миллионщиков В.М. {\em Показатели Ляпунова как функции 


параметра} // Матем. сб. -- 1988. -- Т.\,137. -- \No\,3. -- С.\,364--380. 


\bibitem{4} 


Куратовский К. {\em Топология}. Т.\,1. -- М.: Мир, 1966. -- 594 с. 


\bibitem{5} 


Добронравов В.В. {\em Основы аналитической механики}. -- М.: 


Высш. школа, 1976. -- 264 с. 


\bibitem{6} 


Номидзу К. {\em Группы Ли и дифференциальная геометрия}. -- 


М.: Мир, 1973. -- 188 с. 


\bibitem{7} 


Уиттекер К.Т. {\em Аналитическая динамика}. -- М.--Л.: ОНТИ, 1937. -- 


500 с. 


 


\end{thebibliography} 


 


\vskip 2cm 


 


\post{\it Московский государственный}{\it Поступила} 


\post{\it авиационный институт}{24.03.1999} 


\post{\it $($технический университет$)$}{} 


 


\label{end} 


\end{document} 


